 \subsection{$\mathcal{O}(n)$ 时间复杂度的矩阵-向量乘法}
\label{ssec:jacobian-patt}

PCG 算法涉及频繁的矩阵-向量乘法操作，针对矩阵 $\mathbf{A} = \J^\top\mathbf{H}\J$。我们将证明，在我们的紧凑表示下，将 $\mathbf{H}, \J, \J^\top, \mathbf{A}$ 与向量相乘都可以在 $\mathcal{O}(n)$ 时间内实现。

为简化表述，我们采用如下记号：
\begin{itemize}
  \item $\LowerTri$：块元素均为 $\I$ 的块下三角矩阵，
  \item $\UpperTri$：块元素均为 $\I$ 的块上三角矩阵，
  \item $\Diag$：块对角矩阵（忽略其具体数值），
  \item $\TriDiag$：块三对角矩阵（忽略其具体数值）。
\end{itemize}

\textbf{$\mathbf H$ 的模式：}对于杆仿真，在应用变换 $\mathcal T$ 之前，系统矩阵具有简单的带状结构，采用 $(\mathbf{p}, \mathbf{q})$ 排序时如下：
\begin{equation}
  \label{eq:H-patt}
  \mathbf{H} =
  \begin{bmatrix}
    \mathbf{H}_{\mathbf{p}} & \mathbf{0}              \\
    \mathbf{0}              & \mathbf{H}_{\mathbf{q}}
  \end{bmatrix} =
  \begin{bmatrix}
    \frac{1}{\Delta t^{2}}\mathbf{M}_{\mathbf{p}} & \mathbf{0}                                                               \\
    \mathbf{0}                                    & \frac{1}{\Delta t^{2}} \mathbf{M}_{\mathbf{q}} + \mathbf{K}_{\mathbf{q}}
  \end{bmatrix} =
  \left[
    \begin{array}{@{}c|c@{}}
      \Diag      & \mathbf{0} \\
      \hline
      \mathbf{0} & \TriDiag
    \end{array}
    \right],
\end{equation}
其中 $ \mathbf{H}_{\p} = \FPP{^2E}{\mathbf{p}^2} = \FPP{^2T_{p}}{\mathbf{p}^2} =\mathbf{M}_\p/\Delta t^2$，且 $\mathbf{H}_{\q}= \FPP{^2E}{\mathbf{q}^2} = \FPP{^2T_{q}(\mathbf{q})}{\mathbf{q}^2} + \FPP{^2V(\mathbf{q})}{\mathbf{q}^2} = \mathbf{M}_q /\Delta t^2 + \mathbf{K}_q$。在这种简单的带状结构下，$\mathbf H$ 的矩阵-向量乘法的 $\mathcal{O}(n)$ 时间复杂度自然得到保证。

\textbf{$\mathbf{J}$ 的模式：}我们将雅可比矩阵 $\J$ 分为两块：
\begin{equation}
  \J =
  \begin{bmatrix}
    \J_{\mathbf{p, s}} \\
    \J_{\mathbf{q, s}}
  \end{bmatrix}.
\end{equation}
由于 $\mathbf s=[\p_{\text{c}}^\top, \mathbf{\omega}^\top]^\top$，$\J_{\mathbf{q, s}}$ 具有如下模式：
\begin{equation}
  \J_{\mathbf{q, s}} =
  \begin{bmatrix}
    \frac{\partial \mathbf{q}}{\partial \mathbf{p}_{c}} & \frac{\partial \mathbf{q}}{\partial \mathbf{\omega}}
  \end{bmatrix} =
  \begin{bmatrix}
    \mathbf{0} & \Diag
  \end{bmatrix}.
  \label{eq:Jqs_pattern}
\end{equation}
$\J_{\mathbf{p, s}}$ 的模式相对复杂，如下所示：
\begin{equation}
  \begin{aligned}
    \J_{\mathbf{p, s}} & =
    \begin{bmatrix}
      \I     &                        &                                                 &                        \\
      \I     & \J^{1}_{\mathbf{p, s}} & \multicolumn{2}{c}{\raisebox{1ex}[0pt]{\Huge0}}                          \\
      \vdots & \vdots                 & \ddots                                          &                        \\
      \I     & \J^{1}_{\mathbf{p, s}} & \cdots                                          & \J^{n}_{\mathbf{p, s}} \\
    \end{bmatrix} \\
                       & =
    \begin{bmatrix}
      \I     &        &                                                      \\
      \I     & \I     & \multicolumn{2}{c}{\raisebox{1ex}[0pt]{\Huge0}}      \\
      \vdots & \vdots & \ddots                                          &    \\
      \I     & \I     & \cdots                                          & \I \\
    \end{bmatrix}
    \underbrace{
    \begin{bmatrix}
        \I &                        &        &                        \\
           & \J^{1}_{\mathbf{p, s}} &        &                        \\
           &                        & \ddots &                        \\
           &                        &        & \J^{n}_{\mathbf{p, s}}
      \end{bmatrix}}_{\mathrm{Diag}\left(\{\J^{i}_{\mathbf{p, s}}\}\right)}                                      \\&=
    \begin{bmatrix}
      \LowerTri \cdot \Diag
    \end{bmatrix}.
  \end{aligned}
\end{equation}
显然，我们可以在 $\mathcal{O}(n)$ 时间内将 $\J_{\mathbf{q, s}}$ 与向量相乘。此外，$\J_{\mathbf{p, s}}$ 的矩阵-向量乘法的时间复杂度与 $\LowerTri$ 相同，相当于在 $\mathcal{O}(n)$ 时间内计算向量的前缀和~\cite{harris2007parallel}。因此，我们可以在 $\mathcal{O}(n)$ 时间内实现 $\mathbf J$ 的矩阵-向量乘法。$\J^{\top}$ 具有相同的性质。

最终，我们可以在 $\mathcal{O}(n)$ 时间内将 $\mathbf{A}$ 与向量相乘：
\begin{equation}
  \mathbf{A} = \J^{\top}\mathbf{H}\J =
  \underbrace{
    \left[
      \begin{array}{@{}c|c@{}}
        \Diag \cdot \UpperTri &
        \begin{bmatrix}
          \mathbf{0} \\
          \Diag
        \end{bmatrix}
      \end{array}
      \right]}_{\J^{\top}}
  \underbrace{
    \left[
      \begin{array}{@{}c|c@{}}
        \Diag      & \mathbf{0} \\
        \hline
        \mathbf{0} & \TriDiag
      \end{array}
      \right]}_{\mathbf{H}}
  \underbrace{
    \left[
      \begin{array}{c}
        \LowerTri \cdot \Diag \\
        \hline
        \begin{bmatrix}
          \mathbf{0} & \Diag \\
        \end{bmatrix}
      \end{array}
      \right]}_{\J}.
\end{equation}